\section{Shell-averaged control of the physical zero mode}
\label{sec:zero-shell}

The nonzero Ramanujan--Gauss estimate used later has no content at
additive frequency zero.  We therefore isolate that frequency before
introducing multiplicative characters.  For an off-diagonal row pair
write \(\Delta=m-n\ne0\), and define
\begin{align}
 K^0_{S,T}(m,n)
 :={}&
 \sum_{\substack{u\le S,\ S<v\le T\\
                  (u,mH)=(v,nH)=1}}
 \frac{b_R(u)a(v)}{uv}
 (u,v)\one_{(u,v)\mid\Delta}.
 \label{eq:zero-kernel}
\end{align}
All effective summands are squarefree because both divisor
coefficients are supported on squarefree integers.

\begin{theorem}[Pointwise shell cancellation]
\label{thm:pointwise-zero-shell}
Suppose \(S=X^{s+o(1)}\), \(T\le X^{C_T}\), where \(s,C_T>0\) are
fixed.  For every fixed \(A>0\) and \(0<\kappa<s\), uniformly over the
retained off-diagonal row pairs,
\begin{equation}
 \boxed{
 K^0_{S,T}(m,n)
 \ll_{A,\kappa,C_T,h}
 (\log X)^{-A}+X^{-s+\kappa+o(1)}.}
 \label{eq:pointwise-zero-shell}
\end{equation}
The \(o(1)\) in the second term comes only from the divisor bound for
the nonzero row determinant.
\end{theorem}

\begin{proof}
Put \(g=(u,v)\), \(u=gc\), and \(v=gd\).  Exactness of the gcd and
squarefree support make \(g,c,d\) pairwise coprime.  Reordering the
finite sum gives
\begin{align}
 K^0_{S,T}(m,n)
 ={}&
 \sum_{\substack{g\mid\Delta\\g\ \mathrm{squarefree}\\
                  (g,mnH)=1}}\frac1g
 \sum_{\substack{c\le S/g\\(gc,mH)=1}}
 \frac{b_R(gc)}c
 \notag\\
 &\quad\times
 \sum_{\substack{S/g<d\le T/g\\(d,gcnH)=1}}
 \frac{a(gd)}d.
 \label{eq:zero-gcd-reparametrization}
\end{align}
The modulus in the innermost coprimality condition is exactly
\(gcnH\): the factors \(g,c\) come from \((u,v)=g\), and \(nH\)
comes from target primitivity on the second row.

Split \eqref{eq:zero-gcd-reparametrization} at
\begin{equation}
 G=SX^{-\kappa}.
 \label{eq:g-split}
\end{equation}
If \(g\le G\), then \(S/g\ge X^\kappa\).  Fix \(g,c\) and set
\(M=gcnH\).  Since \(gc\le S\), \(n\asymp Q\), and every scale is a
fixed power of \(X\), there is a fixed \(C_0\) such that
\(M\le X^{C_0}\).  The coefficient identity
\begin{equation}
 a(gd)=\mu(g)a(d)-\mu(g)\log g\,\mu(d)
 \label{eq:a-gd-split}
\end{equation}
therefore turns the inner sum into
\begin{align}
 \mu(g)\{A_M(T/g)-A_M(S/g)\}
 -\mu(g)\log g\{B_M(T/g)-B_M(S/g)\}.
 \label{eq:inner-endpoint-difference}
\end{align}
By \cref{lem:uniform-companion}, the main constants
\(M/\varphi(M)\) in the first difference cancel exactly; both
differences have arbitrary logarithmic decay.  Moreover,
\[
 \frac M{\varphi(M)}\ll\log\log(3X),\qquad
 \sum_{g\mid\Delta}\frac1g\ll\log\log(3X),
\]
and \cref{lem:prefix-mass} bounds the \(c\)-sum by a fixed power of
\(\log X\).  Increasing the calibration exponent in terms of the
desired \(A\) proves
\begin{equation}
 K^{0,\,g\le G}_{S,T}(m,n)\ll_A(\log X)^{-A}.
 \label{eq:small-g-zero}
\end{equation}

For \(g>G\), use absolute values.  The cofactor estimates
\begin{align}
 \sum_{c\le S/g}\frac{|b_R(gc)|}{c}
 &\ll(\log(2X))^4,
 \label{eq:large-g-c-mass}\\
 \sum_{d\le T/g}\frac{|a(gd)|}{d}
 &\ll(\log(2X))^2
 \label{eq:large-g-d-mass}
\end{align}
follow from \cref{lem:prefix-mass} and a harmonic logarithmic sum.
Since \(\Delta\ne0\),
\begin{equation}
 \sum_{\substack{g\mid\Delta\\g>G}}\frac1g
 \le\frac{\tau(|\Delta|)}G
 =X^{-s+\kappa+o(1)}.
 \label{eq:large-g-divisor-saving}
\end{equation}
Combining \eqref{eq:small-g-zero}--
\eqref{eq:large-g-divisor-saving} proves the theorem.
\end{proof}

\subsection{From one row pair to the packet}

Let \(\cZ^0_{S,T}\) denote the contribution of
\eqref{eq:raw-joint-zero} to the joint term in
\eqref{eq:joint-minus-drift}.  Thus
\begin{equation}
 \cZ^0_{S,T}
 =J\widehat F(0)\frac{\varphi(H)}H
 \sum_{m,n}x_my_n\mathfrak m(m,n)K^0_{S,T}(m,n).
 \label{eq:packet-zero-mode}
\end{equation}

\begin{corollary}[Mass-sensitive zero-mode theorem]
\label{cor:mass-sensitive-zero}
Under the hypotheses of \cref{thm:pointwise-zero-shell},
\begin{equation}
 \boxed{
 |\cZ^0_{S,T}|
 \ll_{A,\kappa,F,H}
 J\cZ_1\left\{(\log X)^{-A}
 +X^{-s+\kappa+o(1)}\right\}.}
 \label{eq:mass-sensitive-zero}
\end{equation}
Consequently:
\begin{enumerate}[label=\textup{(\roman*)}]
 \item under only \eqref{eq:row-pair-mass},
 \begin{equation}
  |\cZ^0_{S,T}|
  \ll XQ\left\{X^\epsilon(\log X)^{-A}
  +X^{-s+\kappa+\epsilon+o(1)}\right\};
  \label{eq:soft-mass-zero}
 \end{equation}
 \item under \eqref{eq:polylog-projective-mass}, for every fixed
 \(B>0\),
 \begin{equation}
  |\cZ^0_{S,T}|\ll_{B,K,F,H}XQ(\log X)^{-B};
  \label{eq:polylog-mass-zero}
 \end{equation}
 \item if \(\widehat F(0)=0\), then \(\cZ^0_{S,T}=0\) identically,
 independently of the row mass.
\end{enumerate}
\end{corollary}

\begin{proof}
Insert \eqref{eq:pointwise-zero-shell} into
\eqref{eq:packet-zero-mode}.  Since \(JQ^2\asymp XQ\), the soft
statement follows from \eqref{eq:row-pair-mass}.  It is important that
\(X^\epsilon(\log X)^{-A}\) is retained: for fixed \(A,\epsilon>0\)
it does not tend to zero.  Under
\eqref{eq:polylog-projective-mass}, choose the calibration exponent in
\cref{thm:pointwise-zero-shell} after \(B,K\); the large-\(g\) term
has a fixed power saving because \(\kappa<s\).  The last statement is
immediate from \eqref{eq:packet-zero-mode}.
\end{proof}

\begin{remark}[Why the diagonal mask is structural]
\label{rem:zero-diagonal}
If \(m=n\), then every \(g\) divides \(\Delta=0\), so
\eqref{eq:large-g-divisor-saving} is unavailable.  The row-diagonal
removal is therefore not a cosmetic pruning inherited from an earlier
argument; it is used directly in the zero-mode theorem.
\end{remark}
